\section{Calibrated Magnitude Estimate in the Fourier-Valid Regime}\label{sec:calibration}

\subsection{Atomistic Calibration of $\ell$ and Non-Identifiability of $\gamma$ in FCC Aluminum}\label{subsec:al_calibration}

To assess the physical magnitude of the bulk dispersion and attenuation differences derived in Section~\ref{sec:mode_count}, we evaluate them for face-centered cubic (fcc) Aluminum (Al) at $T_0 = 300\;\mathrm{K}$ using published, independently calibrated atomistic length scales:
\begin{enumerate}[leftmargin=*,align=left]
  \item \textbf{Atomistic 2NN MEAM calibration (Lazar, Agiasofitou, and B{\"o}hlke \cite{Lazar2022}):} Using the Second Nearest-Neighbor Modified Embedded Atom Method (2NN MEAM) interatomic potential for fcc Aluminum (lattice parameter $a = 4.04950\;\mathrm{\AA} = 0.404950\;\mathrm{nm}$), Lazar, Agiasofitou, and B{\"o}hlke \cite{Lazar2022} computed the sixth-order strain-gradient constitutive tensor and extracted the cubic longitudinal characteristic length $\ell_1 = 1.19303\;\mathrm{\AA} = 0.119303\;\mathrm{nm}$ (Table~3 of \cite{Lazar2022}) and the Voigt-averaged isotropic longitudinal characteristic length (Table~6 of \cite{Lazar2022}):
  \begin{equation}\label{eq:lazar_ell}
    \bar{\ell}_1(\mathrm{Al}) = 1.20272\;\mathrm{\AA} = 0.120272\;\mathrm{nm}, \qquad \frac{\bar{\ell}_1}{a} = 0.2970.
  \end{equation}
  \item \textbf{Brillouin zone-edge lattice dispersion match (Eringen \cite{Eringen1972Linear,Eringen2002}):} Matching the 1D Helmholtz dispersion relation $\omega^2(k) = c_0^2 k^2/(1+\ell^2 k^2)$ to the Born--von K{\'a}rm{\'a}n chain at $k = \pi/a$ gives $e_0 = \sqrt{\pi^2 - 4}/(2\pi) \approx 0.39$, or
  \begin{equation}\label{eq:eringen_ell}
    \ell_{\mathrm{Eringen}}(\mathrm{Al}) = 0.39\,a = 0.157931\;\mathrm{nm}.
  \end{equation}
\end{enumerate}

\paragraph{Degeneracy of $(\gamma,\ell)$ at leading order $O(k^2\ell^2)$.}
Expanding the two-phase symbol \eqref{eq:L_symbol} at long wavelengths ($k\ell \ll 1$) gives $\hat{\Lcal}_{\gamma,\ell}(k^2) = 1 - (1-\gamma)k^2\ell^2 + (1-\gamma)k^4\ell^4 - O(k^6\ell^6)$. Matching the $O(k^2)$ dispersion shift to the first strain-gradient dispersion $\omega^2(k) = c_0^2 k^2(1 - \bar{\ell}_1^2 k^2 + O(k^4))$ fixes only the single effective length
\begin{equation}\label{eq:ell_eff}
  \ell_{\mathrm{eff}} := \sqrt{1-\gamma}\,\ell = \bar{\ell}_1 = 0.120272\;\mathrm{nm}.
\end{equation}
As emphasized by Nobili and Pramanik \cite{NobiliPramanik2025b} (Section~2.1), the two-phase local fraction $\gamma \in (0,1)$ has no independent experimental calibration and acts as a singular-perturbation regularization parameter ensuring coercivity on $L^2(\Omega)$ and boundary-value well-posedness.

\subsection{Scale Separation $\chi_{\mathrm{nl}} = \ell/\ell^*$ and Validity Bounds of Fourier Conduction}\label{subsec:scale_separation}

For fcc Aluminum at $T_0 = 300\;\mathrm{K}$ ($\rho = 2700\;\mathrm{kg/m^3}$, $c_0 = 6420\;\mathrm{m/s}$, $c_E = 900\;\mathrm{J/(kg\cdot K)}$, $\kappa = 237\;\mathrm{W/(m\cdot K)}$, dimensionless thermoelastic coupling $\epsilon = 0.0356$), the characteristic thermoelastic frequency $\omega^*$ and thermal length $\ell^*$ are
\begin{equation}\label{eq:al_thermal_scales}
  \omega^* = \frac{\rho c_E c_0^2}{\kappa} = 4.226 \times 10^{11}\;\mathrm{rad/s} \quad (f^* = \frac{\omega^*}{2\pi} = 67.26\;\mathrm{GHz}), \qquad \ell^* = \frac{c_0}{\omega^*} = 15.19\;\mathrm{nm}.
\end{equation}
Consequently, the dimensionless ratio $\chi_{\mathrm{nl}} := \ell/\ell^*$ between the atomistic nonlocal length and the characteristic thermoelastic length is two orders of magnitude smaller than unity:
\begin{equation}\label{eq:chi_nl_values}
  \chi_{\mathrm{nl}} = \frac{\ell}{\ell^*} = \begin{cases}
    7.917 \times 10^{-3} \quad (\chi_{\mathrm{nl}}^2 = 6.268 \times 10^{-5}), & \text{for } \ell = \bar{\ell}_1 = 0.1203\;\mathrm{nm} \;\text{\cite{Lazar2022}}, \\[4pt]
    1.040 \times 10^{-2} \quad (\chi_{\mathrm{nl}}^2 = 1.081 \times 10^{-4}), & \text{for } \ell = 0.39\,a = 0.1579\;\mathrm{nm} \;\text{\cite{Eringen2002}}.
  \end{cases}
\end{equation}

Furthermore, Fourier's heat conduction law $q = -\kappa\partial_x\theta$ is physically defensible only when (i)~the wave period exceeds the microscopic electron--phonon thermal relaxation time $\tau_q \approx 1.5\;\mathrm{ps}$ ($\omega\tau_q \le 0.1 \implies \Omega := \omega/\omega^* \le 0.158$, $f \le 10.6\;\mathrm{GHz}$) and (ii)~the acoustic wavelength $\lambda = 2\pi/k$ exceeds the thermal carrier mean free path $\Lambda_{\mathrm{th}} \approx 15\;\mathrm{nm}$ by an order of magnitude ($\lambda \ge 10\Lambda_{\mathrm{th}} \implies \Omega \le 0.636$, $f \le 42.8\;\mathrm{GHz}$; cf.\ Guyer and Krumhansl \cite{GuyerKrumhansl1966}). At the strict Fourier validity ceiling $\Omega = 0.158$, we have $k\ell = 1.25 \times 10^{-3}$ and $(k\ell)^2 = 1.56 \times 10^{-6}$; even at the Zener thermoelastic peak $\Omega = 1.0$ ($f = 67.26\;\mathrm{GHz}$, $\lambda = 96.3\;\mathrm{nm}$), we have $k\ell \le 1.03 \times 10^{-2}$ and $(k\ell)^2 \le 1.06 \times 10^{-4}$.

\subsection{Quantitative Comparison Across Frequency Regimes}\label{subsec:quantitative_table}

Table~\ref{tab:al_damping} lists the exact dimensionless quasi-elastic attenuation $\mathrm{Im}(k\ell^*)$ and the relative discrepancies $|\alpha_1^{\mathrm{V1}} - \alpha_1^{\mathrm{V2}}|/\alpha_1^{\mathrm{V1}}$ and $|\alpha_1^{\mathrm{Class\,I}} - \alpha_1^{\mathrm{V1}}|/\alpha_1^{\mathrm{V1}}$ for fcc Aluminum at $300\;\mathrm{K}$ ($\gamma = 0.5$) using the atomistic characteristic length $\bar{\ell}_1 = 0.120272\;\mathrm{nm}$ \cite{Lazar2022}.

\begin{table}[htbp]
  \centering
  \caption{Quasi-elastic wave attenuation comparison for fcc Aluminum at $T_0 = 300\;\mathrm{K}$ ($\epsilon = 0.0356$, $\omega^* = 4.226\times 10^{11}\;\mathrm{rad/s}$, $\ell^* = 15.19\;\mathrm{nm}$, $\bar{\ell}_1 = 0.1203\;\mathrm{nm}$ from Lazar et al.\ \cite{Lazar2022}, $\chi_{\mathrm{nl}} = 7.917\times 10^{-3}$, $\gamma = 0.5$).}\label{tab:al_damping}
  \vspace{4pt}
  \begin{tabular}{ccccccp{4.6cm}}
    \toprule
    $\Omega = \omega/\omega^*$ & $f$ (GHz) & $\lambda$ (nm) & $k\ell$ & $\dfrac{|\alpha_1^{\mathrm{V1}} - \alpha_1^{\mathrm{V2}}|}{\alpha_1^{\mathrm{V1}}}$ & $\dfrac{|\alpha_1^{\mathrm{Class\,I}} - \alpha_1^{\mathrm{V1}}|}{\alpha_1^{\mathrm{V1}}}$ & \textbf{Regime Assessment} \\
    \midrule
    $0.001$ & $0.067$ & $97,137$ & $7.78 \times 10^{-6}$ & $6.50 \times 10^{-11}$ & $1.25 \times 10^{-10}$ & Fourier valid ($\omega\tau_q = 6.3\times 10^{-4}$) \\
    $0.010$ & $0.673$ & $9,714$ & $7.78 \times 10^{-5}$ & $5.32 \times 10^{-9}$ & $1.02 \times 10^{-8}$ & Fourier valid ($\omega\tau_q = 6.3\times 10^{-3}$) \\
    $0.100$ & $6.726$ & $971.2$ & $7.78 \times 10^{-4}$ & $\mathbf{5.34 \times 10^{-7}}$ & $\mathbf{1.02 \times 10^{-6}}$ & \textbf{Strict Fourier ceiling} ($\omega\tau_q = 0.063 \le 0.1$) \\
    $0.500$ & $33.63$ & $193.7$ & $3.90 \times 10^{-3}$ & $\mathbf{1.40 \times 10^{-5}}$ & $\mathbf{2.35 \times 10^{-5}}$ & Pre-ballistic transition ($\omega\tau_q = 0.32$) \\
    $1.000$ & $67.26$ & $96.32$ & $7.85 \times 10^{-3}$ & $\mathbf{5.98 \times 10^{-5}}$ & $\mathbf{7.80 \times 10^{-5}}$ & \textbf{Zener peak} ($\omega\tau_q = 0.63$, diff $< 10^{-4}$) \\
    $10.00$ & $672.6$ & $9.53$ & $7.93 \times 10^{-2}$ & $6.28 \times 10^{-3}$ & $4.79 \times 10^{-3}$ & Fourier invalid ($\omega\tau_q = 6.3 \gg 1$, $\lambda < \Lambda_{\mathrm{th}}$) \\
    $100.0$ & $6,726$ & $0.84$ & $8.98 \times 10^{-1}$ & $6.58 \times 10^{-1}$ & $1.01 \times 10^{0}$ & Continuum \& Fourier invalid ($\lambda \approx 2a$) \\
    \bottomrule
  \end{tabular}
\end{table}

Table~\ref{tab:al_damping} establishes a critical physical conclusion that must be stated plainly: \textbf{within the frequency regime where Fourier heat conduction is physically defensible ($\Omega \lesssim 0.16$--$1.0$), the relative difference in bulk thermoelastic attenuation between Variant~V1, Variant~V2, and Class~I in calibrated fcc Aluminum is of order $O(10^{-6}\text{--}10^{-4})$ ($0.0001\%$ to $0.01\%$), and is identical under Eringen's calibration \eqref{eq:eringen_ell} ($\le 1.35 \times 10^{-4}$ at $\Omega = 1$).} An $O(1)$ divergence between the models occurs only when $k\ell \sim 1$ ($\Omega \sim 100$, $f \sim 6.7\;\mathrm{THz}$, $\lambda \sim 0.84\;\mathrm{nm} \approx 2a$), where the wavelength is comparable to two atomic spacings and both Fourier heat conduction ($\omega\tau_q \approx 63$) and continuum elasticity break down. Accordingly, the importance of enforcing Variant~V1 or Variant~V2 in place of Class~I or Class~II is \emph{mathematical and foundational}---guaranteeing compliance with the Second Law of Thermodynamics, correct spatial differential order ($N=4$ vs.\ $N=3$), and well-posed boundary closure on bounded domains---rather than predicting a large measurable bulk ultrasonic damping split in monolithic crystals.
